\documentclass{book}%
\usepackage{graphicx} % pour l'image de la coupe d'Onésime
\usepackage[french]{babel}
\usepackage{fontspec}%
\usepackage{xunicode}%
%\usepackage{yfonts}%
%\usepackage{hyperref}% pour le courriel
\usepackage{amsmath}%
%\usepackage{cancel}% pour \inRas
\usepackage{quoting}%
\usepackage{ulem}% pour souligner en pointillé
\usepackage{comment}%
\usepackage{xcolor}%
\usepackage{graphicx}
%scale = {hscale, vscale}
%\usepackage[scale={0.75,0.8}]{geometry}%
%\setlength{\hoffset}{-18pt}        
%\setlength{\oddsidemargin}{20pt} % Marge gauche sur pages impaires ATTENTION écrit 2 fois


\newfontfamily{\policeFr}{GFS Bodoni}%
%\newfontfamily{\policeLat}{Book Antiqua}%
\newfontfamily{\policeLat}{GFS Neohellenic}%
%\newfontfamily{\policeAn}{Comic Sans MS}%
\newfontfamily{\policeAn}{Courier 10 Pitch}%
%\newfontfamily{\policeAll}{Fraktur}%
%\newfontfamily{\policeAll}{fraktur-bt.ttf}%
\newfontfamily{\policeAll}{Bertholdr Mainzer Fraktur}%
%\newfontfamily{\policeIt}{BodoniFLF}%
\newfontfamily{\policeIt}{GFS Artemisia}%
%\newfontfamily{\policeEsp}{Great Moravia}%
%\newfontfamily\policeGr{GFSBodoni.otf}[BoldFont = GFSBodoniBold.otf, ItalicFont = GFSBodoniIt.otf, BoldItalicFont = GFSBodoniBoldIt.otf]%
%\setmainfont{GFSBodoni.otf}[BoldFont = GFSBodoniBold.otf, ItalicFont = GFSBodoniIt.otf, BoldItalicFont = GFSBodoniBoldIt.otf]
\newfontfamily\policeGr{GFS Bodoni}%
\setmainfont{GFS Bodoni}
%\setdefaultlanguage[variant=ancient]{greek}%
\usepackage{polyglossia}%
\setdefaultlanguage{french}%
%\setotherlanguage[variant=ancient]{greek}%
%\setmainlanguage[variant=ancient]{greek} 
%\setotherlanguage[variant=ancient]{greek} 
%\setmainfont{Libertinus Serif} 
%\newfontfamily{\greekfont}{Liberation Serif}[
%  Script=Greek,
%  NFSSFamily=liberationserif,
%]

\usepackage[maxlevel=3]{csquotes}%
%\MakeOuterQuote{"}% pour les guillemets français
%\usepackage{verse}%


\newcommand{\italique}[1]{\textit{#1}}%
%\newcommand{\citComm}[1]{\begin{quoting}[vskip=0pt] #1 \end{quoting}}% vskip = \topsep + \partopsep par défaut
\newcommand{\espace}[1]{\hspace{#1}}%
\newcommand{\cesure}[0]{ [...] }%
\newcommand{\nieme}[1]{#1\italique{e}}%
\newcommand{\lettreBreve}[1]{$\breve{\mbox{#1}}$}
\newcommand{\synizese}[1]{$\widearc{\mbox{#1}}$}
\newcommand{\inRas}[2]{$\overset{\text{#2}}{\text{\sout{#1}}}$}
\newcommand{\auDessus}[2]{$\overset{\text{#2}}{\text{#1}}$}
%\newcommand{\etc}[0]{κτλ}%
%\newcommand{\cit}[0]{\espace{1cm}} % indent. pour une citation
\newcommand{\expo}[1]{\textsuperscript{#1}}%
\newcommand{\alinea}[0]{ \ \ }%
\newcommand{\Syn}[0]{ Συναγωγή }%
\newcommand{\SynB}[0]{ Συναγωγή\expo{B} }% la deuxième version pour la lettre A
\newcommand{\ClemAlex}{ Clemens Alexandrinus }
\newcommand{\Eustathe}{ Eustathius }
\newcommand{\AelDion}{ Aelius Dionysius }
\newcommand{\EtymMagnum}{ Etymologicum magnum }
\newcommand{\EtymGen}{ Etymologicum genuinum }
\newcommand{\EtymSym}{ Etymologicum Symeonis }
\newcommand{\Apostol}{ Michael Apostolius }
\newcommand{\Phryn}{ Phrynichus \italique{Praep. soph.} }
\newcommand{\Stobee}{ Joannes Stobaeus }
\newcommand{\EpimAlphab}{ Epimerismi alphabetici in Homerum (ed. Cramer) }
\newcommand{\GregCor}{ Gregorius Corinthius }
\newcommand{\Antiatt}{ Antiatticista (ed. Bekker) }
\newcommand{\Harp}{ Harpocration }% Autre édition de référence : Dindorf
\newcommand{\Rpc}[0]{R\expo{pc}}%
\newcommand{\Rac}[0]{R\expo{ac}}%
\newcommand{\Rdeux}[0]{R\expo{2}}%
\newcommand{\Mon}[0]{ M }%
\newcommand{\PsiUn}[0]{PSI-1194}
\newcommand{\PsiDeux}[0]{POxy-3839}
\newcommand{\PsiTrois}[0]{POxy-3840}
\newcommand{\PsiQuatre}[0]{POxy-4935}

\newcommand{\francais}[1]{{\policeFr #1}}%
\newcommand{\latin}[1]{{\policeLat #1}}%
\newcommand{\allemand}[1]{{\policeAll #1}}%
%\newcommand{\allemand}[1]{{\textfrak{#1}}}%
%\newcommand{\grec}[1]{{\policeGr #1}}%
\newcommand{\grec}[1]{#1}%
\newcommand{\anglais}[1]{{\policeAn \scriptsize{#1}}}%
\newcommand{\italien}[1]{{\policeIt #1}}%
\newcommand{\espagnol}[1]{{\policeEsp #1}}%

\newcommand{\Nauck}[0]{N\expo{2}}%
\newcommand*{\auteur}[1]{\textsc{#1}*}%
\newcommand{\cf}[0]{\italique{cf.}\xspace}%
\newcommand{\chap}[1]{\^{#1}}%

\newlength{\larg}%
\newcommand{\decalage}[1]{%
\settowidth{\larg}{#1}%
\espace{\larg}%
}
\newcommand{\decalageN}[1]{%
\settowidth{\larg}{#1}%
\espace{-\larg}%
}

\usepackage{xcolor}
\newcommand{\fantome}{{\color{white}v}}

\newcommand{\lecon}[1]{\textbf{#1}}
%\newcommand{\douteux}[1]{\dotuline{#1}} % pour les lettres douteuses : un point en-dessous
\newcommand{\douteux}[1]{$\underset{\cdot}{\text{#1}}$}
\newcommand{\difference}[1]{\dotuline{#1}}

\newcommand{\marqueur}[1]{\textsuperscript{#1}}%

\newcommand{\eff}[1]{%
{\color{orange}#1}%
}%

\newcommand{\LSJ}[0]{LSJ}%



\newcommand{\Index}[1]{\section*{#1}%
\addcontentsline{toc}{section}{#1}}
\newcommand{\note}[1]{\footnote{#1}}
\newcommand{\biblio}[2]{\hspace*{-1.5cm}\textbf{[#1]}\hspace*{1.5cm}#2\\

}
\newcommand{\Auteur}[1]{\hspace*{-1.5cm}\textbf{#1} \\}
\newcommand{\livre}[1]{\hspace*{0.5cm}#1\\}
\newcommand{\passage}[3]{\hspace*{1.0cm} #1 \hspace*{1.2cm}\decalageN{#1} #2 \dotfill \hfill	\textbf{#3}\\} % passage Corrigé ou expliqué
%\newcommand{\passageC}[3]{\hspace*{2.5cm} #1 \hspace*{1cm} #2 \dotfill \hfill	\textbf{#3}\\} % passage Corrigé ou expliqué

\pagenumbering{alph}


\usepackage{minipage-marginpar}
\usepackage{marginnote}
\usepackage[marginparwidth=3cm]{geometry}% pour les marges
%\usepackage{geometry}
\usepackage{multicol}

\begin{document} 
\renewcommand{\contentsname}{Index}
\title{Index}
\author{}\date{}
\maketitle
\setcounter{page}{0}


\tableofcontents
%\setcounter{page}{0}
\cleardoublepage
%\setcounter{page}{1}


%
%\Index{Index des auteurs cités}
%dsvdsvs
%\begin{multicols}{2}
%texte à placer sur N colonnes
%	
%fnrene
%
%rzbqze
%
%gzqerg
%dvqezzzzzzzzzzzzzzzzzzzzzzz
%
%
%zrgzg
%\end{multicols}

\Index{Index des allusions proverbiales}
(Une liste des proverbes cités par Platon -- qu'il faudrait certainement compléter -- a été dressée par Jowett, \italique{The dialogues of Plato}, V, pp. 497-498.)


\Auteur{Aristophane}
\livre{\italique{Acharniens}}
\passage{258}{περιτράγῃ τὰ χρυσία}{894}
\livre{\italique{Nuées}}
\passage{55}{σπαθᾶν [τὴν κρόκην]}{738}
\passage{394}{`βροντὴ' καὶ `πορδή' ὁμοίω}{49-50 n. 2}
\livre{\italique{Guêpes}}
\passage{170}{αὐτοῖσι τοῖς κανθηλίοις}{819-23}
\livre{\italique{Paix}}
\passage{793-5}{ἔφασκε [$=$ Καρκίνος] τὸ δρᾶμα [$=$ Μῦες] γαλῆν τῆς ἑσπέρας ἀπάγξαι}{558-9}
\passage{1253}{συρμαίαν μετρεῖν $\sim$ κύματα μετρεῖν}{855-7}
\livre{\italique{Thesmophories}}
\passage{19}{μήτ' ἀκούω μήθ' ὁρῶ}{19}
\passage{21}{οἷόν τέ που ’στιν αἱ σοφαὶ ξυνουσίαι}{21}
\passage{24}{προσμάθοιμι χωλὸς εἶναι τὼ σκέλει}{22-4}
\passage{94}{ἡμέτερος ὁ πυραμοῦς}{94}
\passage{140}{τίς δαὶ κατόπτρου καὶ ξίφους κοινωνία}{140}
\passage{168}{ὁ Φιλοκλέης αἰσχρὸς $=$ κορυδός}{168-70}
\passage{241-2}{ὕδωρ ὦ γείτονες πρὶν ἀντιλαβέσθαι πρωκτὸν ἓτερον τῆς φλογός}{241-2}
\passage{247}{μὴ φροντίσῃς· ἕτερος σπογγιεῖ (cf. \italique{Grenouilles} 487 ?)}{247 n. 1}
\passage{248}{οἰμώξετ' ἄρ' εἴ τις τὸν ἐμὸν πρωκτὸν πλυνεῖ}{248}
\passage{394}{τὰς οὐδὲν ὑγιές. Cf. Érasme n° 738, \italique{nihil sanum}.}{394}
\passage{455-6}{ἄγρια δρᾷ κακά ἅτ' ἐν ἀγρίοισι τοῖς λαχάνοις αὐτὸς τραφείς}{455-6}
\passage{489}{ἐχομένη τῆς δάφνης}{489}
\passage{527}{ἀλλ' ἅπαν γένοιτ' ἂν ἤδη}{527}
\passage{529-30}{ὑπὸ λίθῳ παντί ἀθρεῖν}{528-30}
\passage{530}{δάκῃ ῥητωρ $\sim$ δῆγμα συκοφάντου}{528-30}
\passage{531-2}{οὐδὲν κάκιον φύσει τῶν γυναικῶν πλὴν ἄρ' εἰ γυναῖκες}{531-2}
\passage{559}{ἔπειτα τὴν γαλῆν φαμεν}{558-9}
\passage{568}{λαβὲ θοἰμάτιον Φιλίστη}{568}
\passage{633}{σκάφιον Ξένυλλ' ᾔτησεν· οὐ γὰρ ἦν ἀμίς}{633}
\passage{702}{ἅπαντα τόλμης κἀναισχυντίας}{702}
\passage{738}{κακὸν δὲ καὶ τοῖς σκευαρίοις καὶ τῇ κρόκῃ}{738}
\passage{773}{ἀντὶ τῶν πλατῶν $\sim$ ἄλλο γένος κώπης ?}{773-4}
\passage{822}{αὐτῶ κανῷ}{819-23}
\passage{894}{ἐπὶ κλοπῇ τοῦ χρυσίου}{894}
\passage{928}{αὕτη μὲν ἡ μήρινθος οὐδὲν ἔσπασεν}{928}
\passage{938}{καίπερ ἀποθανουμένῳ $\sim$ Ἀττικὸς ὑπέχει τὴν χεῖρα ἀποθνήσκων}{936-7}
\passage{940}{πρὸς τῇ σανίδι δεῖν $=$ δῆσαι πρὸς σανίδι (Hesychius δ 926)}{939-42}
\passage{949}{Παύσων σέβεται καὶ νηστεύει $\sim$ Παύσωνος πτωχότερος}{949-52}
\passage{1079}{ἐς κόρακας}{1079}
\passage{1192}{ὠς γλυκερὸ τὸ γλῶσσ', ὤσπερ Ἀττικὸς μέλις $\sim$ μελίγλωσσος}{1192}
\livre{\italique{Grenouilles}}
\passage{172}{βούλει σκευάρι' εἰς Ἅιδου φέρειν}{738}
\passage{560}{αὐτοῖς τοῖς ταλάροις}{819-23}





\Index{Index des gloses d'Hésychius utiles pour l'établissement et la bonne compréhension du texte}
(Je ne recense ici que les gloses d'Hésychius qui ont servi au rétablissement de l\allemand{Urtext} ou à la bonne compréhension d'un passage encore mal expliqué. Il n'y a qu'à jeter un œil sur les \latin{testimonia} pour constater combien les commentaires les plus érudits d'Aristophane ont scrupuleusement été dépouillés et intégrés au lexique d'Hésychius.)\\
\Auteur{Aristophane}
\livre{\italique{Cavaliers}}
\passage{20}{α 6416 : ἀπόκυνον : μάζα [ ... ] πρὸς ἀναίρεσιν κυνῶν}{455-6}
\livre{\italique{Oiseaux}}
\passage{1126}{δ 2227 : ἐλύμνιαι· δοκοὶ ὀρόφιναι}{1069}
\livre{\italique{Thesmophories}}
\passage{50}{π 3188 : πραῖνοι· πρηνίζει, καταστρέφει}{49-50}
\passage{80}{τ 1440 : τρίτη Θεσμοφορίων· ζητεῖται πῶς ἅμα μὲν λέγει τρίτην Θεσμοφορίων εἶναι, ἅμα δὲ μέσην, τεττάρων οὐσῶν ἡμερῶν}{80}
\passage{100}{μ 1904 : Μύρμηκος ἀτραπούς· [ ... ] ἀπὸ ἥρωος Μύρμηκος ὀνομαζομένη}{100}
\passage{163}{δ 1805 : δίκανα· ποικίλα ἱμάτια}{161-3}
\passage{291}{θ 50 : θάλλικα· σάκκου εἶδος}{291}
\passage{704}{ε 3599 : ἐξαρῆξαι· ἐκφορῆσαι}{704}
\passage{743}{τ 1388 : τρικότυλος οἶνος· οὗ τρεῖς κοτύλαι ὀβολοῦ πωλοῦνται}{743}
\passage{781}{τ 1162 : τορεῖν· τορῆσαι, τρῆσαι, τεμεῖν}{781}
\passage{853}{ε 1995 : ἑλένη· λαμπάς, δετή}{853}
\passage{894}{χ 792 : χρυσίον· τὸ τῶν παιδίων αἰδοῖον}{894 n.5}
\passage{910}{λ 171 : λαίφια· ῥάκη}{910}
\passage{940}{δ 926 : δῆσαι πρὸς σανίδι· Ἀττικῶς· βασανίσαι ἐν τῷ δεσμωτηρίῳ}{939-42}
\passage{945}{κ 4170, 4270 : κρόκαλον· τὸ πανοῦργον παιδίον (κρόκαλ' $\sim$ κροκώθ' ?)}{945}
\passage{1041}{ψ 135 : ψεφαυγοῦς· σκοτεινῆς (cf Schol. ad v.1054)}{1041}
\Auteur{Euripide}
\livre{\italique{Fragments} (ed. Kannicht)}
\passage{228}{κ 468 : †κάλλης (immo καλλήϊον)· καρπίμου}{855-7}
\passage{228}{γ 66 : Γαιών· ποταμός, παρ' Ἕλλησι Νεῖλος}{855-7}



\clearpage



\Index{Index des corrections proposées}
\Auteur{Aristophane}
\livre{\italique{Cavaliers}}
\passage{20}{ἀπόκινον codices omnes : ἀπόκυνον Bibi ap. Hesych. α 6416}{455-6}
\livre{\italique{Lysistrata}}
\passage{104;114}{Les deux répliques sont prononcées par la même femme}{643}
\livre{\italique{Thesmophories}}
\passage{1}{χελιδὼν R : Χελιδὼν Bibi ap. Clements}{1-10}
\passage{31}{Ἀγάθων R : ἀγαθῶν ; Bibi}{30-3}
\passage{50}{πράμος R : πραῖνοι Bibi ap. Hesych. π 3188}{49-50}
\passage{71-4}{(Μνη.) (71) et (Εὐρ.) (72-4) R : (Μνη.) (71- 4) Bibi ap. Schol.}{70-4}
\passage{163}{διεκίνων R : διεκίνουν Suid. : δίκανον Bibi ap. Hesych. δ 1805}{161-3}
\passage{217}{διδόναι γ’ αὐτὸν R : διδόναι γ’ ἑαυτὸν Bibi auctore Willems}{217}
\passage{291}{πρὸς θάληκον R : Ποσθαλίσκον Dindorf : πρὸς θάλικα Bibi}{291}
\passage{446}{ἐν Κύπρῳ R : ἐγκύμων Bibi}{446}
\passage{557}{τὸν σῖτον R, Schol., Suid. σ 512 : τὸν κύσθον Bibi, cf. v. 1114}{556-7}
\passage{637-40}{(Μνη.) R : (Μνη.) (637-8) et (Μίκα) (639-40) Bibi}{637-40}
\passage{638}{χάλᾳ R : χάλα Grynaeus : χαλᾷ Bibi}{637-40}
\passage{702}{ἔργα R : πλέα Suid. : δοῦλα cod. Laur. : διπλὰ Bibi}{702}
\passage{704}{ἐξάρξω R : ἐξαράξω Bentley : ἐξαρήξω Bibi}{704}
\passage{710}{ἥκεις R : ἥκεις Suid. : ἥκων Willems : ἥκειν Bibi ap. Phot. η 88}{710-2}
\passage{711}{οὐ λέξεις R, Suid. : οὐκ ἔξει. Bibi ap. Schol.}{711-2}
\passage{721}{ἀθέοις ἔργοις R : ξὺν ἀθέοις ἔργοις Bibi ap. Taillardat, cf. v. 716}{721-5 n.3}
\passage{743}{τρικότυλον ; ἢ πῶς. --- εἰπέ μοι· τί μ' ἠργάσω ; (ap. POxy-3839)}{743}
\passage{781}{τὸ ῥῶ R : τορῶ Bibi ap. Hesych. τ 1162}{781}
\passage{822-3}{ταὐτὸ κανοῦν \inRas{οἱ καλαθίσκοι}{} καὶ τὸ σκιάδειον}{819-23}
\passage{839}{ὦ πόλεις R : ὦ πόλις Gelenius : ωὖπολι Bibi}{839}
\passage{853}{Ἑλένη R : Ἑλένην Bonini : ἑλένην Bibi, ap. Hesych. ε 1995}{853}
\passage{910}{ἀφύων R : ἰφύων Schol., Suid. : ἀμφίων Grégoire : λαιφίων Bibi}{910}
\passage{1031}{κημὸν ἕστηκ' ἔχουσ' R : ἕστηκ' ἔχουσα κημόν, cf. schol. \italique{Vesp.} 674c}{1031}
\passage{1041}{φεύγουσαν R : ψεφαυγοῦς Bibi, cf Schol. ad v.1054, Hesych. ψ 135}{1041}
\passage{1071}{μέρος ἐξέλαχον; R : μόρος ἔξελ' ἄχος. Bibi}{1071}
\passage{1074}{στωμυλωμένη R : στωμυλιομενής Bibi}{1073-5}
\passage{1078}{παῦσαι R : παΰσαι Bibi}{1078}
\passage{1079}{βάλλ' ἐς κόρακας R : βάλλε σύ κόρακας Bibi}{1079}
\passage{1093}{(Εὐρ.) ποῖ ποῖ πέτει om. R : suppl. Bibi}{1093}
\passage{1151}{θεμιτὸν R : θέμις Burges : θεμιτ' Bothe : θεμιτεύετον Bibi}{1150-3}
\passage{Σ 1174}{τὸ ἁβρῶς βαδίζειν R : τὸ ἄκρως βαδίζειν Bibi ap. Zonaras}{1174}
\passage{1197}{συμβήνην R, Bibi : συβήνην Schol., Suid., edd. omn.}{1196-217}
\livre{\italique{Grenouilles}}
\passage{19-20}{οὐκ ἐρεῖ codd. : οὐ κε ῥεῖ Bibi}{647-8 n.3}
\livre{\italique{Fragments} (ed. Kassel-Austin)}
\passage{192.96}{[τὸ]\douteux{ν} ναυτικὸν \douteux{ἄ}[ργυρον}{839}
\Auteur{Diogenianus}
\passage{IV 78}{ταὐτὸ τῇ codd. : ταὐτότητι Bibi}{489}
\Auteur{Euripide}
\livre{\italique{Fragments} (ed. Kannicht)}
\passage{228}{κάλλιστον ἐκ γαίας : καλλήϊον Γαιᾶ Bibi ap Hesych. κ 468, γ 66}{855-7}
\Auteur{Eustathe de Thessalonique}
\livre{\italique{Commentaire à l'}Iliade}
\passage{976.13}{λέγεται δὲ φησίν codd. : λέγεται δὲ σφῇσιν Bibi}{643}
\Auteur{Hipponax}
\livre{\italique{Fragments} (ed. Masson / Degani)}
\passage{30 / 42}{οὔ μοι δικαίως μοιχὸς edd. : οὔ μοι δίκᾳ ὡς μοιχὸς Bibi}{894}
\Auteur{Sappho}
\livre{\italique{Fragments} (ed. Reinach-Puech / Éva-Maria Voigt)}
\passage{70/125}{αὐτὰ κόραις στεφαναπλόκεο}{400-1}





\clearpage



\Index{Index des passages davantage explicités.}
{\centering }
\Auteur{Aristophane}
\livre{\italique{Cavaliers}}
\passage{765}{[Κλέων] βέλτιστος ἀνὴρ μετὰ Λυσικλέα καὶ Κύνναν καὶ Σαλαβακχώ}{805}
\livre{\italique{Nuées}}
\passage{133}{κόρακες $=$ \italique{marteaux de porte}}{1079}
\passage{394}{`βροντὴ' καὶ `πορδή' ὁμοίω}{49-50 n. 2}
\passage{508}{εἰς Τροφωνίου}{570}
\livre{\italique{Guêpes}}
\passage{98-9}{Δῆμον καλόν ... κημὸς καλός}{1031}
\passage{853}{τοὺς καδίσκους $=$ \italique{les muselières}}{1031}
\livre{\italique{Paix}}
\passage{793-5}{[Καρκίνος] ἔφασκε τὸ δρᾶμα [$=$ Μῦες] γαλῆν τῆς ἑσπέρας ἀπάγξαι}{558-9}
\passage{1252-3}{πώλει [τὰ κράνη] τοῖς Αἰγυπτίοις ... ἐπιτήδεια συρμαίαν μετρεῖν}{855-7}
\livre{\italique{Oiseaux}}
\passage{1126}{Aller se faire pendre εἰς Ἐλύμνιον}{1069}
\livre{\italique{Lysistrata}}
\passage{111-4}{τοὔγκυκλον τουτὶ καταθεῖσαν ἐκπιεῖν αὐθημερόν}{643}
\passage{681}{αὐχήν}{647-8 n. 2}
\livre{\italique{Thesmophories}}
\passage{1-10}{ὦ Ζεῦ ... οὐχ ἅ γ᾽ ἂν ἀκούειν δέῃ.}{1-10}
\passage{80}{ἐπεὶ τρίτη ’στὶ Θεσμοφορίων ἡ μέση}{80}
\passage{98}{Κυρήνην δ' ὁρῶ}{98}
\passage{100}{μύρμηκος ἀτραπούς}{100}
\passage{101-29}{Chant d'Agathon}{101-29}
\passage{273}{τὴν Ὑποκράτους ξυνοικίαν}{273}
\passage{336-7}{ἢ 'πικηρυκεύεται Εὐριπίδῃ Μήδοις ἐπὶ βλάβῃ τινὶ.}{336-7}
\passage{346}{ἢ καὶ δέχεται προδιδοῦσ' ἑταίρα τὸν φίλον}{346}
\passage{455-6}{ἄγρια δρᾷ κακά ἅτ' ἐν ἀγρίοισι τοῖς λαχάνοις αὐτὸς τραφείς}{455-6}
\passage{487}{ἐγὼ δὲ καταχέασα τοῦ στροφέως ὕδωρ}{487}
\passage{500}{ὑπ' αὐγὰς οἷόν ἐστιν ἐγκεκαλυμμένον}{500}
\passage{529-30}{ὑπὸ λίθῳ γὰρ παντί που χρὴ μὴ δάκῃ ῥητωρ ἀθρεῖν}{528-30}
\passage{556-7}{στλεγγίδας λαβοῦσαι ἔπειτα σιφωνίζομεν τὸν κύσθον}{556-7}
\passage{570}{τὸν σησαμοῦνθ' ὃν κατέφαγες, τοῦτον χεσεῖν ποιήσω}{570}
\passage{605}{Κλεώνυμος / ῥίψασπις}{605}
\passage{633}{σκάφιον Ξένυλλ' ᾔτησεν· οὐ γὰρ ἦν ἀμίς}{633}
\passage{643}{ἔγκυκλον}{643}
\passage{647}{ἰσθμόν}{647-8}
\passage{747}{σχεδὸν τοσοῦτον χὤσον ἐκ Διονυσίων.}{747}
\passage{754}{δός \{μοι\} τὸ σφαγεῖον}{754}
\passage{760-1}{τίς ἐξεκόρησέ σε; τίς τὴν ἀγαπητὴν παῖδά σοὐξηράσατο;}{760-1}
\passage{782}{χώρει, χώρει}{782}
\passage{782-4}{ποίαν αὔλακα ... ταχέως χρή}{782-4}
\passage{804}{Ναυσιμάχης ἥττων ἐστὶν Χαρμῖνος}{804}
\passage{805}{Κλεοφῶν χείρων πάντως δήπου Σαλαβακχοῦς}{805}
\passage{811}{κατὰ πεντήκοντα τάλαντα, cf. Hérodote (VI 132-136)}{811-2}
\passage{826}{ἐκ τῶν οἴκων $=$ τὰ πατρῷα}{826}
\passage{841-2}{Ὑπέρβολος ... Λάμαχος}{841-2}
\passage{855-7}{Νείλου ῥοαί ὃς Αἰγύπτου πέδον λευκῆς νοτίζει μελανοσυρμαῖον λεών}{855-7}
\passage{870}{μὴ ψεῦσον ὦ Ζεῦ τῆς ἐπιούσης ἐλπίδος}{870}
\passage{898}{εἰ μὴ Κρίτυλλά γ' Ἀντιθέου Γαργηττόθεν}{898}
\passage{910}{ἐγὼ δὲ Μενελέῳ σ' ὅσα γ' ἐκ τῶν λαιφίων}{910}
\passage{915}{ἄπαγέ μ' ἄπαγ' ἄπαγέ με}{913-5}
\passage{936-7}{ὦ πρύτανι πρὸς τῆς δεξιᾶς, ἥνπερ φιλεῖς κοίλην προτείνειν}{936-7}
\passage{939-42}{γυμνὸν ἀποδύσαντά με / πρὸς τῇ σανίδι δεῖν / ἐν κροκωτοῖς μίτραις}{939-42}
\passage{949}{Παύσων σέβεται καὶ νηστεύει}{949-52}
\passage{1015-55}{Les phrases musicales semblent aller par paires $\sim$ échos.}{1015-55}
\passage{1059}{Ἠχὼ λόγων ἀντῳδὸς ἐπικοκκάστρια}{1059}
\passage{1136-59}{Chant ibycéen du chœur}{1136-59}
\passage{1174}{πρῶτον μὲν οὖν δίελθε κἀνακόλπασον}{1174}
\passage{1175}{Περσικόν... distinguer la danse de l'air de flûte}{1175}
\passage{1192}{ὠς γλυκερὸ τὸ γλῶσσ ̓, ὤσπερ Ἀττικὸς μέλις $\sim$ μελίγλωσσος}{1192}
\passage{1197}{ἀλλὰ τὸ συμβήνην λαβέ ... ὀρτῶς δὲ συβήνη ’στί· καταβηνῆσι γάρ}{1196-217}
\livre{\italique{Grenouilles}}
\passage{19-20}{ἄρ' ὁ τράχηλος οὑτοσί, ὅτι θλίβεται μέν, τὸ δὲ γέλοιον οὐκ ἐρεῖ}{647-8 n.3}
\livre{\italique{Assemblée des femmes}}
\passage{535-8}{ἐπιβαλοῦσα τοὔγκυκλον μόνον οὐ στεφανώσασ’ οὐδ’ ἐπιθεῖσα λήκυθον}{643}
\passage{741}{ὄρθριος νόμος}{643 n. 4}
\livre{\italique{Fragments} (ed. Kassel-Austin)}
\passage{192.96-9}{ἐκδανείζει τὸν ναυτικὸν ἄργυρον ἐπὶ τῷ πέμπτῳ μέρει τοὺς τόκους}{839}
\Auteur{Athénée}
\livre{\italique{Deipnosophistes}}
\passage{I 4d}{ἀμίδα δότω τις ἢ πλακοῦντά τις δότω}{570}
\Auteur{Euripide}
\livre{\italique{Fragments} (ed. Kannicht)}
\passage{228}{Νείλου λιπὼν ὕδωρ ὃς πληροῦται ῥοὰς  ἡνίκ’ ἂν τακῇ χιὼν}{855-7}



























%\Index{Index général}
\end{document}